\subsection{Construção da rede dirigida e ponderada}
\label{sec:construcao-rede}

A rede tem 67 nós, um por atividade da MIP, identificados pelos códigos do IBGE. Cada nó representa uma origem emissora nas ligações de saída e o setor de um produto final nas ligações de entrada, conforme a representação de \textcite{huang2019}. Em cada nó de destino, a atribuição reúne as emissões associadas aos diferentes usos finais dos produtos da atividade, como consumo das famílias e exportações. A matriz de adjacência ponderada desta rede é definida como:
\begin{equation}
W_{ij}=\begin{cases}
F_{ij}, & i\ne j,\\
0, & i=j.
\end{cases}
\label{eq:rede-sem-diagonal}
\end{equation}
Cada $W_{ij}>0$ define uma aresta $i\rightarrow j$, cujo peso, em Gg de CO\textsubscript{2}, indica a emissão em $i$ atribuída ao produto final de $j$. As duas direções de um par são distintas. Como os pesos incorporam requisitos diretos e indiretos, concatenar arestas não representa uma sequência de etapas produtivas.

A diagonal $F_{ii}$ reúne emissões de $i$ atribuídas à demanda final de seus próprios produtos, incluindo requisitos indiretos. Ela permanece em $F$ e nas contas $c$ e $u$, mas é retirada de $W$, cujo peso total corresponde apenas às atribuições entre atividades distintas.

\subsection{Caracterização da rede}
\label{sec:volume-concentracao}

Para caracterizar a presença de relações intersetoriais, calculam-se o grau de cada atividade e a densidade da rede.

\paragraph{Grau.}
O grau de saída $k_i^{\mathrm{out}}$ conta os setores de produtos finais aos quais uma atividade atribui emissões; o grau de entrada $k_i^{\mathrm{in}}$ conta as origens emissoras que compõem a atribuição aos produtos finais dessa atividade:
\begin{equation}
\begin{aligned}
k_i^{\mathrm{out}}&=\left|\{j\ne i:W_{ij}>0\}\right|,\\
k_i^{\mathrm{in}}&=\left|\{j\ne i:W_{ji}>0\}\right|.
\end{aligned}
\label{eq:graus-rede}
\end{equation}
As barras indicam o número de elementos de cada conjunto.

\paragraph{Densidade.}
A densidade é a proporção de ligações possíveis que estão presentes na rede. Como há 67 atividades e as ligações de uma atividade consigo mesma são excluídas, existem $67\times66$ ligações dirigidas possíveis:
\begin{equation}
d=\frac{\sum_i k_i^{\mathrm{out}}}{67\times66}.
\label{eq:densidade-rede}
\end{equation}
A densidade varia de zero, quando não há ligações, a um, quando todas as ligações possíveis estão presentes.

Para examinar como as emissões se distribuem entre atividades emissoras e setores de produtos finais, analisam-se as forças de saída e de entrada e a concentração das relações representadas em W. A análise conjunta das forças e concentrações permite distinguir atividades com volumes semelhantes de emissões intersetoriais, mas com diferentes distribuições entre origens e destinos. \parencite{huang2019,serrano2009}.

\paragraph{Forças.}
A força é a soma dos pesos das ligações de um nó em uma dada direção. Neste trabalho, corresponde ao volume intersetorial de emissões, em Gg de CO\textsubscript{2}.
As forças de saída e de entrada de cada atividade são calculadas pela soma dos pesos das respectivas ligações:
\begin{equation}
\begin{aligned}
s_i^{\mathrm{out}}&=\sum_{j\ne i}W_{ij}=c_i-F_{ii},\\
s_i^{\mathrm{in}}&=\sum_{j\ne i}W_{ji}=u_i-F_{ii}.
\end{aligned}
\label{eq:forcas-externas}
\end{equation}
A força de saída $s_i^{\mathrm{out}}$ mede as emissões próprias de $i$ atribuídas aos produtos finais de outras atividades; a força de entrada $s_i^{\mathrm{in}}$ mede as emissões de outras atividades atribuídas aos produtos finais de $i$. A diferença entre força de entrada e força de saída reproduz o saldo entre as contas: $s_i^{\mathrm{in}}-s_i^{\mathrm{out}}=u_i-c_i$.

\paragraph{Participações.}
A participação expressa quanto cada ligação representa na força da atividade:
\begin{equation}
p_{ij}^{\mathrm{out}}=\frac{W_{ij}}{s_i^{\mathrm{out}}},\qquad
p_{ji}^{\mathrm{in}}=\frac{W_{ji}}{s_i^{\mathrm{in}}},\quad j\ne i.
\label{eq:participacoes-externas}
\end{equation}
As participações somam um em cada perfil de saída ou de entrada.

\paragraph{Concentrações}
O índice de concentração $Y$\footnote{Na terminologia de redes, essa medida também é denominada disparidade dos pesos.} mede em que medida as participações se concentram em poucas relações:
\begin{equation}
Y_i^{\mathrm{out}}=\sum_{j\ne i}(p_{ij}^{\mathrm{out}})^2,\qquad
Y_i^{\mathrm{in}}=\sum_{j\ne i}(p_{ji}^{\mathrm{in}})^2.
\label{eq:disparidade}
\end{equation}
Quanto maior $Y$, maior a concentração. Para $k$ relações de igual peso, $Y=1/k$; quando todo o peso está associado a uma única relação, $Y=1$.

\paragraph{Número efetivo de contrapartes.}
O número efetivo de contrapartes, $D_2$, expressa quantas atividades com participações iguais produziriam a concentração observada. É calculado pelo inverso do índice de concentração:
\[
D_{2,i}^{\mathrm{out}}=1/Y_i^{\mathrm{out}},\qquad
D_{2,i}^{\mathrm{in}}=1/Y_i^{\mathrm{in}}.
\]
Para $k$ pesos iguais, $D_2=k$; quando uma relação domina, $D_2$ aproxima-se de um. Valores menores de $D_2$ indicam maior concentração. 

\subsection{Extração e avaliação do backbone}
\label{sec:backbone}

O filtro de disparidade é um método que permite reduzir a densidade da rede ao selecionar ligações com participações excepcionalmente altas nos pesos de uma origem ou de um destino \parencite{serrano2009,huang2019}. A sub-rede formada pelas ligações selecionadas é denominada backbone e preserva seus pesos originais, sem alterar as contas ambientais completas.

Na rede $W$, $k$ é o número de ligações de saída ou de entrada de um nó, conforme a direção analisada. O modelo nulo representa a distribuição aleatória da força correspondente entre essas ligações dividindo o intervalo de zero a um em $k$ parcelas, por meio de $k-1$ pontos escolhidos uniformemente ao acaso. Cada parcela representa a participação de uma ligação na força do nó.

\paragraph{Critério de seleção.}
Para $k>1$, a probabilidade de uma parcela atingir ou superar $p$ é $(1-p)^{k-1}$. Para a aresta $i\rightarrow j$, as caudas nas perspectivas de saída e entrada são
\begin{equation}
\begin{aligned}
a_{ij}^{\mathrm{out}}&=\left(1-\frac{W_{ij}}{s_i^{\mathrm{out}}}\right)^{k_i^{\mathrm{out}}-1},\\
a_{ij}^{\mathrm{in}}&=\left(1-\frac{W_{ij}}{s_j^{\mathrm{in}}}\right)^{k_j^{\mathrm{in}}-1}.
\end{aligned}
\label{eq:caudas-filtro}
\end{equation}
A aresta é selecionada se $a_{ij}^{\mathrm{out}}<\alpha$ \textbf{ou} $a_{ij}^{\mathrm{in}}<\alpha$, registrando-se a extremidade responsável pela seleção. A saída destaca um destino nas atribuições da origem; a entrada destaca uma origem na composição do destino.\footnote{As caudas são calculadas apenas para graus pertinentes maiores que um. Se uma extremidade tem grau um, vale o teste da outra; se ambas têm grau um, a aresta é preservada por convenção de manutenção da relação exclusiva, inclusive em díades isoladas (par de nós que possui ligação entre si, mas nenhuma ligação com outros nós da rede). Essa extensão dos casos degenerados (grau zero ou um) do suplemento de \textcite{serrano2009} é registrada separadamente da seleção estatística.} Adota-se $\alpha=0{,}05$ como limiar convencional de referência, com $0{,}01$ e $0{,}10$ para sensibilidade \parencite{serrano2009}. 

Para cada valor de $\alpha$, o filtro é aplicado uma vez a $W$, gerando a matriz filtrada $B_\alpha$. Essa matriz mantém todos os nós e preserva os pesos originais das arestas que satisfazem o critério de seleção, atribuindo zero às demais, sem renormalização ou corte absoluto adicional.

\paragraph{Avaliação da seleção.}
Para cada $\alpha$, compara-se a rede filtrada $B_\alpha$ com a rede completa $W$ para avaliar a redução do número de ligações e a parcela das emissões intersetoriais preservada. A fração de arestas retidas é a razão entre o número de arestas positivas em $B_\alpha$ e em $W$. A cobertura compara a soma dos pesos retidos com a soma dos pesos originais, tanto para a rede inteira quanto para cada atividade e direção:
\begin{equation}
\begin{aligned}
r_\alpha&=\frac{\sum_{ij}(B_\alpha)_{ij}}{\sum_{ij}W_{ij}},\\
r_{\alpha,i}^{\mathrm{out}}&=\frac{\sum_j(B_\alpha)_{ij}}{s_i^{\mathrm{out}}},\qquad
r_{\alpha,i}^{\mathrm{in}}=\frac{\sum_j(B_\alpha)_{ji}}{s_i^{\mathrm{in}}}.
\end{aligned}
\label{eq:cobertura-backbone}
\end{equation}
A cobertura global $r_\alpha$ divide o peso total de $B_\alpha$ pelo peso total de $W$. Para a atividade $i$, $r_{\alpha,i}^{\mathrm{out}}$ divide a soma da linha $i$ de $B_\alpha$ pela força de saída em $W$, enquanto $r_{\alpha,i}^{\mathrm{in}}$ divide a soma da coluna $i$ de $B_\alpha$ pela força de entrada em $W$. As razões são calculadas para denominadores positivos e excluem a diagonal. Essa comparação permite verificar se uma cobertura global elevada convive com perdas expressivas nas relações de determinadas atividades.

A conectividade é avaliada comparando-se os graus, o número de nós isolados e os componentes fracamente conectados de $B_\alpha$ e $W$. Estes últimos são os grupos de nós conectados quando se desconsidera a direção das arestas.

Para avaliar a sensibilidade ao parâmetro, comparam-se os backbones de $\alpha=0{,}01$, $0{,}05$ e $0{,}10$ quanto às arestas selecionadas, às coberturas e às posições das atividades por grau. O aninhamento dos conjuntos decorre da regra e não constitui evidência independente de robustez. As caudas do modelo nulo não medem incerteza ambiental, e a seleção não estabelece prioridade de mitigação.

Para interpretar os perfis, as contrapartes são ordenadas por peso em $W$ completa. Destacam-se as cinco maiores e a participação conjunta das demais, sem modificar a rede. A leitura dos casos combina essas participações, as contas completas e os motivos de retenção; a sensibilidade verifica também a permanência das principais relações nos três parâmetros.
